
\begin{resumo}
Este artigo compara seis algoritmos de classificação aplicados à predição de renda no \textit{Adult Income Dataset}. O objetivo é prever se a renda anual de um indivíduo excede 50 mil dólares a partir de dados censitários. Foram avaliados Árvore de Decisão, KNN, Naive Bayes, Regressão Logística, LinearSVC e MLP. A metodologia utilizou pré-processamento, otimização de hiperparâmetros e validação cruzada estratificada. Os resultados, avaliados por Acurácia, Precisão, Revocação e F1-Score, evidenciaram desempenhos distintos e reforçaram a importância do tratamento de dados heterogêneos e desbalanceados.
\end{resumo}

\begin{abstract}
This paper compares six classification algorithms applied to income prediction using the \textit{Adult Income Dataset}. The goal is to predict whether an individual's annual income exceeds \$50,000 based on census data. Decision Tree, KNN, Naive Bayes, Logistic Regression, LinearSVC, and MLP were evaluated. The methodology used preprocessing, hyperparameter optimization, and stratified cross-validation. The results, assessed through Accuracy, Precision, Recall, and F1-Score, showed different performance profiles and reinforced the importance of handling heterogeneous and imbalanced data.
\end{abstract}


\clearpage